\documentclass[11pt,a4paper]{article}

\usepackage[margin=0.82in]{geometry}
\usepackage[T1]{fontenc}
\usepackage[utf8]{inputenc}
\usepackage{lmodern}
\usepackage[table]{xcolor}
\usepackage{booktabs}
\usepackage{longtable}
\usepackage{tabularx}
\usepackage{array}
\usepackage{enumitem}
\usepackage{listings}
\usepackage[most]{tcolorbox}
\usepackage{titlesec}
\usepackage{fancyhdr}
\usepackage{xurl}
\usepackage[colorlinks=true,linkcolor=nsblue,urlcolor=nsteal,citecolor=nsblue]{hyperref}

\definecolor{nsblue}{HTML}{1F6FB2}
\definecolor{nsteal}{HTML}{187F79}
\definecolor{nsgreen}{HTML}{247A4B}
\definecolor{nsamber}{HTML}{B57500}
\definecolor{nsred}{HTML}{B8322A}
\definecolor{nsgrey}{HTML}{343A40}
\definecolor{nslight}{HTML}{EEF5F7}
\definecolor{codebg}{HTML}{F6F8FA}

\setlength{\parindent}{0pt}
\setlength{\parskip}{0.45em}
\setlength{\headheight}{14pt}
\setcounter{tocdepth}{2}
\setlist[itemize]{leftmargin=1.45em,itemsep=2pt,topsep=3pt}
\setlist[enumerate]{leftmargin=1.65em,itemsep=3pt,topsep=3pt}
\newcolumntype{Y}{>{\raggedright\arraybackslash}X}
\newcolumntype{P}[1]{>{\raggedright\arraybackslash}p{#1}}

\titleformat{\section}{\sffamily\Large\bfseries\color{nsblue}}{\thesection}{0.55em}{}
  [{\color{nsblue!25}\titlerule[0.8pt]}]
\titleformat{\subsection}{\sffamily\large\bfseries\color{nsgrey}}{\thesubsection}{0.5em}{}
\titleformat{\subsubsection}{\sffamily\normalsize\bfseries\color{nsteal}}{}{0pt}{}

\pagestyle{fancy}
\fancyhf{}
\fancyhead[L]{\small\sffamily\color{nsgrey}NeuroShield Software-First MVP Task Plan}
\fancyhead[R]{\small\sffamily\color{nsgrey}\thepage}
\renewcommand{\headrule}{\vskip1pt{\color{nsblue!40}\hrule height 0.7pt}}

\lstset{
  basicstyle=\ttfamily\small,
  backgroundcolor=\color{codebg},
  frame=single,
  rulecolor=\color{black!15},
  breaklines=true,
  columns=fullflexible,
  showstringspaces=false,
  upquote=true,
  xleftmargin=4pt,
  xrightmargin=4pt,
  aboveskip=5pt,
  belowskip=5pt
}

\newtcolorbox{taskbox}[2][]{
  breakable,
  enhanced,
  colback=white,
  colframe=nsblue,
  colbacktitle=nsblue,
  coltitle=white,
  fonttitle=\sffamily\bfseries,
  title={#2},
  boxrule=0.8pt,
  arc=2pt,
  left=8pt,right=8pt,top=7pt,bottom=7pt,
  #1
}

\newtcolorbox{donebox}{
  breakable,
  colback=nsgreen!6,
  colframe=nsgreen,
  title={Definition of done},
  fonttitle=\sffamily\bfseries,
  boxrule=0.9pt,
  arc=2pt
}

\newtcolorbox{warningbox}[1]{
  breakable,
  colback=nsamber!8,
  colframe=nsamber,
  colbacktitle=nsamber!8,
  coltitle=nsamber!65!black,
  fonttitle=\sffamily\bfseries,
  title={#1},
  boxrule=0.9pt,
  arc=2pt
}

\newtcolorbox{plainbox}[1]{
  breakable,
  colback=nslight,
  colframe=nsblue!55,
  title={#1},
  fonttitle=\sffamily\bfseries,
  boxrule=0.8pt,
  arc=2pt
}

\newcommand{\code}{\path}
\newcommand{\green}{\textbf{\textcolor{nsgreen}{GREEN}}}
\newcommand{\amberstatus}{\textbf{\textcolor{nsamber}{AMBER}}}
\newcommand{\redstatus}{\textbf{\textcolor{nsred}{RED}}}

\begin{document}
\thispagestyle{empty}

\begin{tcolorbox}[colback=nsblue,colframe=nsblue,arc=3pt,boxrule=0pt,
  left=16pt,right=16pt,top=16pt,bottom=16pt]
{\sffamily\color{white}\fontsize{27}{31}\selectfont\bfseries NeuroShield}\par
\vspace{4pt}
{\sffamily\color{white}\LARGE Software-First MVP Task Plan}\par
\vspace{8pt}
{\sffamily\color{white!90}\large Self-contained tasks for building the MVP with public datasets and synthetic live streams before hardware}
\end{tcolorbox}

\vspace{8pt}
\textbf{Project file context:} \code{manika/initial_research.tex}\hfill
\textbf{Working assumption:} build software first, connect hardware later

\begin{plainbox}{MVP strategy}
The MVP is a software-first research prototype. First, build a complete laptop app that works with
public datasets and a generated/replayed sensor stream. Only after the replay demo is stable should
hardware be connected. When hardware arrives, it must replace the stream source only; the feature
pipeline, model, backend, dashboard, logging, explanations, and tests should already work.
\end{plainbox}

\begin{warningbox}{Honesty boundary}
NeuroShield is not a clinical device. The MVP detects patterns associated with public stress-proxy
datasets and local simulated streams. It must not claim to predict panic attacks, diagnose anxiety,
diagnose heat illness, or detect burnout. Public datasets used here contain proxy labels such as
baseline, lab stress, task stress, activity, and self-report.
\end{warningbox}

\tableofcontents
\newpage

% ============================================================================
\section{MVP Outcome}

\subsection{What must exist at the end of the software-first MVP}

\begin{longtable}{P{0.25\linewidth}P{0.66\linewidth}}
\toprule
\textbf{Area} & \textbf{Required MVP behavior} \\
\midrule
Dataset ingestion & WESAD loads correctly first. At least one external dataset loader or documented skip decision exists. \\
Feature pipeline & One versioned feature extractor converts windows into heart, EDA, temperature, and motion features. \\
Model & M1 rest-vs-stress-proxy model trains on WESAD with subject-wise evaluation and saves a versioned artifact. \\
Synthetic/replay source & A software stream emits the same raw event schema that later hardware will emit. \\
Personalization & A baseline profile is computed from quiet windows and used to z-score features. \\
Live runtime & Replay/synthetic stream flows through features, baseline, motion gate, model, status, and explanations. \\
Dashboard & Local dashboard shows connection, quality, \green/\amberstatus/\redstatus, motion pause, reasons, and history. \\
Tests and evidence & Dataset smoke tests, feature tests, model metrics, backend tests, dashboard build, and a recorded replay demo exist. \\
\bottomrule
\end{longtable}

\subsection{What is deliberately postponed}

Hardware wiring, firmware, LiPo power, Bluetooth phone app, on-device TinyML, permanent enclosure,
clinical claims, panic prediction, real burnout prediction, and volunteer stress experiments are
postponed until the software replay gate passes.

\subsection{Directory structure to build}

\begin{lstlisting}
neuroshield/
  README.md
  pyproject.toml
  uv.lock
  package.json
  data/
    raw/                 # not committed
    external/            # downloaded public datasets; not committed
    interim/             # converted parquet/npz; not committed
    fixtures/            # tiny synthetic fixtures; committed
  docs/
    contracts.md
    datasets.md
    model_card_m1.md
    software_acceptance.md
  src/neuroshield/
    data/
    features/
    models/
    runtime/
    api/
  app/                   # dashboard
  tests/
  artifacts/
    metrics/
    plots/
    models/
    demo/
\end{lstlisting}

% ============================================================================
\section{Dataset Download and Contents}

Store raw datasets outside Git. The recommended root is \code{data/external/}. Create one folder per
dataset and keep a \code{SOURCE.txt} file inside each folder with the URL, download date, license or
terms note, file names, sizes, and checksum if available.

\begin{warningbox}{Do not download everything on day one if bandwidth is limited}
For the MVP, WESAD is mandatory. PPG-DaLiA is strongly recommended for motion handling. Stress-Predict
or Nurse Stress is useful for external validation. SWELL-KW and Fitbit are optional concept modules
and can be deferred.
\end{warningbox}

\subsection{Dataset priority map}

\begin{longtable}{P{0.17\linewidth}P{0.16\linewidth}P{0.16\linewidth}P{0.42\linewidth}}
\toprule
\textbf{Dataset} & \textbf{Priority} & \textbf{Used for} & \textbf{What to expect} \\
\midrule
WESAD & Must have & M1 stress-proxy model & 15 subjects; wrist Empatica E4 BVP, EDA, temperature, ACC; chest signals; baseline/stress/amusement labels. \\
PPG-DaLiA & Strongly recommended & Motion gate and pulse-quality tests & 15 subjects; wrist PPG and accelerometer during daily activities; chest ECG ground truth for heart rate. Large download. \\
Stress-Predict & Recommended & External sanity check & 35 volunteers; Empatica E4 raw folders and processed CSV; baseline/stress labels across Stroop, interview, hyperventilation. \\
Nurse Stress & Optional validation & Naturalistic stress test & 15 nurses; roughly one week of E4 data; sparse self-report stress surveys; messy real-world data and empty files may appear. \\
SWELL-KW & Optional M2 & Cognitive-load concept & Office-task stress dataset; request/download process may take time; use HRV/EDA only if access is granted. \\
Fitbit Fitness Tracker & Optional trends & BurnoutWatch-style trend demo & Kaggle dataset with activity, heart-rate, and sleep CSVs; no stress or burnout labels. \\
\bottomrule
\end{longtable}

\subsection{WESAD: primary training dataset}

\textbf{Official source:}
\url{https://archive.ics.uci.edu/dataset/465/wesad+wearable+stress+and+affect+detection}

\textbf{What it contains.}
\begin{itemize}
\item 15 subjects in a lab stress-affect protocol.
\item Wrist Empatica E4 channels: BVP at 64 Hz, EDA at 4 Hz, body temperature at 4 Hz, and 3-axis acceleration at 32 Hz.
\item Chest RespiBAN channels: ECG, EDA, EMG, respiration, temperature, and acceleration at 700 Hz.
\item Labels for baseline/neutral, stress, amusement, and other protocol periods.
\item Subject pickle files and a README from the linked original dataset package.
\end{itemize}

\textbf{How to download.}
\begin{enumerate}
\item Open the UCI WESAD page above.
\item Click \textbf{Dataset Home Page}. UCI links to the original University of Siegen WESAD page.
\item Read and follow the dataset acknowledgment/license instructions on the original page.
\item Download the WESAD zip/package from the original dataset page.
\item Save the archive as \code{data/external/wesad/WESAD.zip} or similar.
\item Extract so the expected layout contains subject folders such as \code{S2}, \code{S3}, etc.
\item Create \code{data/external/wesad/SOURCE.txt} with source URL, date, and citation.
\end{enumerate}

\textbf{Command-line option after manual download.}
\begin{lstlisting}
mkdir data\external\wesad
# Move the downloaded WESAD archive into data\external\wesad
# Then extract with Windows Explorer, 7-Zip, or:
tar -xf data\external\wesad\WESAD.zip -C data\external\wesad
\end{lstlisting}

\textbf{Use in MVP.}
Use only wrist E4 signals first because they match the future glove signals most closely. Train M1 on
baseline versus stress. Ignore amusement for the first binary classifier, or keep it only for a later
three-class experiment.

\subsection{PPG-DaLiA: motion and pulse-quality dataset}

\textbf{Official source:}
\url{https://archive.ics.uci.edu/dataset/495/ppg+dalia}

\textbf{What it contains.}
\begin{itemize}
\item 15 subjects performing daily-life activities.
\item Wrist Empatica E4 BVP/PPG, EDA, temperature, and acceleration.
\item Chest RespiBAN ECG, respiration, and acceleration.
\item ECG-derived ground truth heart rate for testing pulse extraction under motion.
\item UCI lists \code{data.zip} at about 2.7 GB plus a README PDF.
\end{itemize}

\textbf{How to download.}
\begin{enumerate}
\item Open the UCI PPG-DaLiA page above.
\item Click \textbf{Download} for \code{data.zip}; also download \code{readme.pdf}.
\item Save them in \code{data/external/ppg_dalia/}.
\item Extract \code{data.zip}; verify subject folders/files are visible.
\item Record source URL, date, file size, and license note in \code{data/external/ppg_dalia/SOURCE.txt}.
\end{enumerate}

\textbf{Command-line option.}
\begin{lstlisting}
mkdir data\external\ppg_dalia
# Browser download is safest because the file is large.
# After download:
tar -xf data\external\ppg_dalia\data.zip -C data\external\ppg_dalia
\end{lstlisting}

\textbf{Use in MVP.}
Do not train the stress classifier on PPG-DaLiA. Use it to decide when the live app should abstain
with \code{motion_paused} or \code{poor_signal}, and to test whether pulse-derived heart-rate features
stay reasonable during movement.

\subsection{Stress-Predict Dataset: external lab validation}

\textbf{Official source:}
\url{https://github.com/italha-d/Stress-Predict-Dataset}

\textbf{What it contains.}
\begin{itemize}
\item 35 healthy volunteers.
\item Stress tasks: Stroop colour-word test, interview session, and hyperventilation, with baseline/relaxation periods.
\item Raw Empatica E4 folders per participant containing \code{ACC.csv}, \code{BVP.csv}, \code{EDA.csv}, \code{HR.csv}, \code{IBI.csv}, \code{TEMP.csv}, \code{info.txt}, and task tags.
\item Processed files including a combined HR/respiration binary-label CSV, questionnaire scores, and task time logs.
\end{itemize}

\textbf{How to download.}
\begin{enumerate}
\item Open the GitHub repository above.
\item Click \textbf{Code} then \textbf{Download ZIP}, or use Git if installed.
\item Save/extract into \code{data/external/stress_predict/}.
\item Confirm both \code{Raw_data/} and \code{Processed_data/} exist.
\item Record repository URL, commit hash if using Git, download date, and license in \code{SOURCE.txt}.
\end{enumerate}

\textbf{Command-line option.}
\begin{lstlisting}
git clone https://github.com/italha-d/Stress-Predict-Dataset.git data\external\stress_predict
cd data\external\stress_predict
git rev-parse HEAD
\end{lstlisting}

\textbf{Use in MVP.}
Use this after WESAD training to see whether the model behaves sensibly on another lab stress
protocol. It is also useful for a software-only stream fixture because the raw E4 CSVs are easy to
replay.

\subsection{Nurse Stress Dataset: naturalistic validation}

\textbf{Official source:}
\url{https://datadryad.org/dataset/doi:10.5061/dryad.5hqbzkh6f}

\textbf{What it contains.}
\begin{itemize}
\item Data from 15 female nurses working hospital shifts.
\item About 1,250 hours collected in two sessions during 2020.
\item Empatica E4 signal files: \code{EDA.csv}, \code{HR.csv}, \code{TEMP.csv}, \code{IBI.csv}, \code{BVP.csv}, and \code{ACC.csv}.
\item \code{SurveyResults.xlsx} with reported stress level and context columns.
\item Some raw CSV files may be empty because of device capture failures; loaders must handle this.
\end{itemize}

\textbf{How to download.}
\begin{enumerate}
\item Open the Dryad page above.
\item Download \code{Stress_dataset.zip} and \code{SurveyResults.xlsx}, or click \textbf{Download full dataset}.
\item Save both under \code{data/external/nurse_stress/}.
\item Extract the zip; confirm participant folders and raw signal CSVs are present.
\item Record DOI, version date, file sizes, and citation in \code{SOURCE.txt}.
\end{enumerate}

\textbf{Use in MVP.}
Use only after the WESAD pipeline is stable. Expect noisy, sparse labels and missing files. Treat it
as a robustness check, not as the primary proof of accuracy.

\subsection{SWELL-KW: optional cognitive-load module}

\textbf{Project page:}
\url{http://cs.ru.nl/~skoldijk/SWELL-KW/Dataset.html}

\textbf{What it contains.}
\begin{itemize}
\item Knowledge-worker office tasks with stressors such as time pressure and email interruption.
\item Physiological signals relevant to HRV and EDA, plus richer modalities that the glove will not reproduce.
\item Questionnaire/task-load information for stress/cognitive-load conditions.
\end{itemize}

\textbf{How to download.}
\begin{enumerate}
\item Open the SWELL-KW dataset page.
\item Follow the page's access/request instructions. This dataset may require a request rather than immediate direct download.
\item If access is granted, save files under \code{data/external/swell_kw/}.
\item Record the access date, granted terms, file names, and citation in \code{SOURCE.txt}.
\end{enumerate}

\textbf{Use in MVP.}
Optional. Use it only for M2 FocusGuard/cognitive-load experiments. The main MVP does not depend on it.

\subsection{Fitbit Fitness Tracker Data: optional trends dataset}

\textbf{Source:}
\url{https://www.kaggle.com/datasets/arashnic/fitbit}

\textbf{What it contains.}
\begin{itemize}
\item Consumer-style activity, heart-rate, and sleep CSVs.
\item Longitudinal daily/minute-level behavior useful for dashboard trend logic.
\item No labelled stress, burnout, or clinical outcomes.
\end{itemize}

\textbf{How to download.}
\begin{enumerate}
\item Create or sign into a Kaggle account.
\item Open \code{https://www.kaggle.com/settings} and create an API token. This downloads \code{kaggle.json}.
\item Place \code{kaggle.json} in \code{\%USERPROFILE\%\.kaggle\kaggle.json} on Windows.
\item Install the Kaggle CLI in the project environment.
\item Run the download command below.
\item Unzip into \code{data/external/fitbit/} and record source/date in \code{SOURCE.txt}.
\end{enumerate}

\begin{lstlisting}
uv add kaggle
uv run kaggle datasets download -d arashnic/fitbit -p data\external\fitbit --unzip
\end{lstlisting}

\textbf{Use in MVP.}
Optional. Use it to build a History/Trends screen that demonstrates rolling averages and recovery
patterns. Do not call it burnout prediction.

% ============================================================================
\section{Software Tasks}

Each task below is written so a builder can complete it without reading later tasks. Complete tasks in
order unless a task explicitly says it is optional.

\begin{taskbox}{T1 -- Create the software repo and scope files}
\textbf{Goal.} Create a project folder that makes the software-first assumption explicit.

\textbf{Steps.}
\begin{enumerate}
\item Create the directory structure from Section 1.3.
\item Add \code{README.md} with the MVP statement: ``software replay first, hardware source later.''
\item Add \code{docs/product_scope.md} listing required release-1 behavior and postponed hardware work.
\item Add \code{docs/no_clinical_claims.md} stating the proxy-label limitation.
\item Add \code{.gitignore} entries for \code{data/raw/}, \code{data/external/}, \code{data/interim/}, \code{artifacts/models/*.joblib}, and local session logs.
\end{enumerate}

\begin{donebox}
The repo can be opened by a new person and they can answer: what is being built now, what is postponed,
where datasets go, and why hardware is not required for the first demo.
\end{donebox}
\end{taskbox}

\begin{taskbox}{T2 -- Install and lock the Python environment}
\textbf{Goal.} Make dataset loading, feature extraction, training, API serving, and tests reproducible.

\textbf{Steps.}
\begin{enumerate}
\item Install \code{uv} if it is not installed.
\item Create \code{pyproject.toml}.
\item Add dependencies: \code{numpy}, \code{pandas}, \code{scipy}, \code{scikit-learn}, \code{neurokit2}, \code{matplotlib}, \code{seaborn}, \code{joblib}, \code{pydantic}, \code{fastapi}, \code{uvicorn}, \code{pyserial}, \code{pytest}, \code{ruff}, and \code{mypy}.
\item Run \code{uv sync}.
\item Add a smoke command that prints package versions.
\end{enumerate}

\begin{lstlisting}
uv init
uv add numpy pandas scipy scikit-learn neurokit2 matplotlib seaborn joblib pydantic fastapi uvicorn pyserial
uv add --dev pytest ruff mypy
uv sync
uv run python -c "import pandas, sklearn, neurokit2, fastapi; print('software env ok')"
\end{lstlisting}

\begin{donebox}
\code{uv.lock} exists, the smoke command passes, and a clean checkout can recreate the same environment
with \code{uv sync}.
\end{donebox}
\end{taskbox}

\begin{taskbox}{T3 -- Write the raw event contract before any streaming code}
\textbf{Goal.} Define the stream format that both generated data and future hardware will use.

\textbf{Steps.}
\begin{enumerate}
\item Create \code{docs/contracts.md}.
\item Define one newline-delimited JSON event per sample or health update.
\item Required common fields: \code{schema_version}, \code{type}, \code{source}, \code{session_id}, \code{seq}, \code{t_us}, and \code{ok}.
\item Define event types: \code{ppg}, \code{eda}, \code{temp}, \code{imu}, and \code{health}.
\item Define units: PPG raw counts, EDA relative level or microsiemens when source supports it, temperature Celsius, acceleration m/s\textsuperscript{2}.
\item Add examples for each event type.
\end{enumerate}

\begin{lstlisting}
{"schema_version":"neuroshield.hw.v1","type":"eda","source":"synthetic",
 "session_id":"demo-001","seq":144,"t_us":4500000,"eda_level":0.42,"ok":true}
\end{lstlisting}

\begin{donebox}
\code{docs/contracts.md} has enough detail that a synthetic generator and future firmware can produce
the same schema without guessing field names or units.
\end{donebox}
\end{taskbox}

\begin{taskbox}{T4 -- Download WESAD and prove one subject loads}
\textbf{Goal.} Get the mandatory training dataset onto disk and inspect it before writing generic code.

\textbf{Steps.}
\begin{enumerate}
\item Follow the WESAD download instructions in Section 2.3.
\item Create \code{src/neuroshield/data/wesad_loader.py}.
\item Implement a function \code{load_wesad_subject(subject_id)}.
\item Return wrist BVP, EDA, TEMP, ACC, labels, sample rates, and subject ID.
\item Plot one subject's BVP, EDA, TEMP, ACC magnitude, and labels for a 10-minute slice.
\item Save the plot to \code{artifacts/plots/wesad_subject_smoke.png}.
\end{enumerate}

\textbf{Expected obstacles.} WESAD files are commonly pickled Python objects. Use \code{pickle.load(..., encoding='latin1')} if needed. Confirm the exact keys by printing one subject dictionary before building the final loader.

\begin{donebox}
One subject loads, the plot shows plausible physiological signals, label segments are visible, and
the loader returns sample rates instead of assuming one shared rate.
\end{donebox}
\end{taskbox}

\begin{taskbox}{T5 -- Convert WESAD into a canonical signal bundle}
\textbf{Goal.} Hide dataset-specific structure behind one stable object used by all later code.

\textbf{Steps.}
\begin{enumerate}
\item Create \code{src/neuroshield/data/bundle.py}.
\item Define a \code{SignalBundle} Pydantic/dataclass object with channels, time arrays, sample rates, labels, subject ID, dataset name, and warnings.
\item Convert WESAD wrist signals into this bundle.
\item Preserve native sample rates: BVP 64 Hz, EDA/TEMP 4 Hz, ACC 32 Hz.
\item Write tests for subject ID, channel names, monotonic time, label presence, and non-empty arrays.
\end{enumerate}

\begin{donebox}
\code{pytest tests/test_wesad_loader.py} passes and no downstream task needs to know WESAD's internal pickle structure.
\end{donebox}
\end{taskbox}

\begin{taskbox}{T6 -- Build the first feature extractor}
\textbf{Goal.} Convert each time window into model-ready features.

\textbf{Steps.}
\begin{enumerate}
\item Create \code{src/neuroshield/features/extract.py}.
\item Use 60-second windows with 30-second overlap for training.
\item Compute feature columns:
\begin{itemize}
\item \code{hr_mean_bpm}
\item \code{ibi_sd_ms}
\item \code{ibi_rmssd_ms}
\item \code{ppg_quality}
\item \code{eda_level}
\item \code{eda_slope}
\item \code{eda_response_count}
\item \code{eda_response_mean_amp}
\item \code{temp_mean_c}
\item \code{temp_slope_c_per_min}
\item \code{motion_dynamic_rms}
\item \code{motion_dynamic_p95}
\item \code{valid_fraction}
\end{itemize}
\item Use NeuroKit2 for PPG/EDA processing where possible.
\item Add a \code{feature_version = "features-v1"} constant.
\item Save one subject's feature table to \code{data/interim/wesad_features/S2.parquet}.
\end{enumerate}

\begin{donebox}
Feature extraction produces a table with fixed column order, no silent all-null columns, one row per
valid window, labels attached, and a test that catches accidental column changes.
\end{donebox}
\end{taskbox}

\begin{taskbox}{T7 -- Label WESAD windows for the M1 binary classifier}
\textbf{Goal.} Create clean baseline-vs-stress labels without data leakage.

\textbf{Steps.}
\begin{enumerate}
\item Map WESAD baseline windows to \code{0}.
\item Map WESAD stress windows to \code{1}.
\item Exclude amusement and transition/undefined windows from M1.
\item Drop windows with too little valid signal coverage.
\item Save a label-count table per subject to \code{artifacts/metrics/wesad_label_counts.csv}.
\end{enumerate}

\begin{donebox}
Every training row has exactly one binary label, subject IDs are preserved, excluded windows are counted,
and class balance is visible before model training.
\end{donebox}
\end{taskbox}

\begin{taskbox}{T8 -- Train M1 with leave-one-subject-out evaluation}
\textbf{Goal.} Produce the headline software model and honest metrics.

\textbf{Steps.}
\begin{enumerate}
\item Create \code{src/neuroshield/models/train_m1.py}.
\item Use \code{LeaveOneGroupOut} with subject ID as the group.
\item Train a scikit-learn pipeline with imputation, scaling if needed, and either logistic regression or random forest.
\item Compare against a dummy classifier.
\item Report balanced accuracy, macro-F1, confusion matrix, and per-subject metrics.
\item Save predictions to \code{artifacts/metrics/m1_loso_predictions.csv}.
\item Save summary metrics to \code{artifacts/metrics/m1_loso_summary.json}.
\end{enumerate}

\begin{donebox}
The model is evaluated subject-wise, a dummy baseline is included, metrics are saved, and no subject's
windows appear in both train and test folds.
\end{donebox}
\end{taskbox}

\begin{taskbox}{T9 -- Save a versioned M1 model artifact}
\textbf{Goal.} Freeze one model for the replay demo.

\textbf{Steps.}
\begin{enumerate}
\item Train a final M1 model on all eligible WESAD windows.
\item Save it as \code{artifacts/models/m1_wesad_features_v1.joblib}.
\item Create \code{artifacts/models/m1_wesad_features_v1_manifest.json}.
\item Manifest must include model version, feature version, feature column order, training dataset, code commit if available, metrics file path, threshold policy, and checksum.
\item Add model-loading code that refuses incompatible feature versions.
\end{enumerate}

\begin{donebox}
The model can be loaded by version, the feature order is checked before inference, and a mismatch fails clearly instead of producing a misleading prediction.
\end{donebox}
\end{taskbox}

\begin{taskbox}{T10 -- Create the synthetic stream generator}
\textbf{Goal.} Generate live-looking data before hardware exists.

\textbf{Steps.}
\begin{enumerate}
\item Create \code{src/neuroshield/runtime/synthetic_source.py}.
\item Generate event streams for phases: quiet baseline, mild stress-like rise, motion burst, recovery, and sensor fault.
\item Emit the same raw event schema from \code{docs/contracts.md}.
\item Include controllable seed, duration, sample rates, and phase schedule.
\item Write a CLI command to save a fixture:
\end{enumerate}

\begin{lstlisting}
uv run python -m neuroshield.runtime.synthetic_source ^
  --out data\fixtures\calm_motion_stress.ndjson ^
  --duration-sec 600 --seed 7
\end{lstlisting}

\begin{donebox}
A committed fixture under \code{data/fixtures/} can drive the whole app without internet, datasets, or hardware.
\end{donebox}
\end{taskbox}

\begin{taskbox}{T11 -- Implement replay source and immutable raw logs}
\textbf{Goal.} Treat generated files and future hardware streams as interchangeable sources.

\textbf{Steps.}
\begin{enumerate}
\item Create \code{src/neuroshield/runtime/replay_source.py}.
\item Read NDJSON events in timestamp order and optionally play them at real speed or accelerated speed.
\item Validate every event against Pydantic schemas.
\item Preserve invalid lines in a raw log with parse errors rather than silently dropping them.
\item Add counters for valid events, invalid events, missing fields, unknown schema versions, and stale periods.
\end{enumerate}

\begin{donebox}
Replay of \code{calm_motion_stress.ndjson} produces the same typed events every run, and malformed fixtures are rejected with stable error messages.
\end{donebox}
\end{taskbox}

\begin{taskbox}{T12 -- Compute personal baseline profiles}
\textbf{Goal.} Make predictions personalized even before hardware.

\textbf{Steps.}
\begin{enumerate}
\item Create \code{src/neuroshield/runtime/baseline.py}.
\item Define a baseline profile JSON schema with feature means, feature standard deviations, accepted seconds, source, feature version, and creation time.
\item From a quiet replay segment, compute baseline mean and standard deviation for every feature.
\item Protect against division by zero using a minimum standard deviation policy.
\item Add a function that transforms a feature row into z-scored features.
\end{enumerate}

\begin{donebox}
A baseline JSON can be created from a quiet fixture, loaded later, rejected if feature versions mismatch,
and used to z-score all model input features.
\end{donebox}
\end{taskbox}

\begin{taskbox}{T13 -- Add motion and signal-quality abstention}
\textbf{Goal.} Prevent the model from guessing when data is unreliable.

\textbf{Steps.}
\begin{enumerate}
\item Define thresholds for \code{motion_dynamic_rms}, \code{motion_dynamic_p95}, \code{ppg_quality}, and \code{valid_fraction}.
\item If motion is too high, output \code{motion_paused}.
\item If signal coverage or PPG quality is too low, output \code{poor_signal}.
\item Keep abstention separate from \green/\amberstatus/\redstatus predictions.
\item Test with the synthetic motion burst.
\end{enumerate}

\begin{donebox}
The replay fixture visibly pauses inference during motion and resumes after the configured quiet period.
No low-quality window receives a confident stress status.
\end{donebox}
\end{taskbox}

\begin{taskbox}{T14 -- Convert model probabilities into live states}
\textbf{Goal.} Turn M1 output into stable product states.

\textbf{Steps.}
\begin{enumerate}
\item Define states: \code{waiting}, \code{calibrating}, \code{green}, \code{amber}, \code{red}, \code{motion_paused}, \code{poor_signal}, \code{stale}, and \code{error}.
\item Define threshold policy: example \green{} below 0.45, \amberstatus{} from 0.45 to 0.70, \redstatus{} above 0.70.
\item Add persistence/hysteresis so one spike does not immediately trigger red.
\item Save every status with timestamp, model version, feature version, probability, quality, and reasons.
\end{enumerate}

\begin{donebox}
A deterministic replay produces a predictable sequence of states: waiting, green, motion paused,
recovery, amber/red test segment, recovery.
\end{donebox}
\end{taskbox}

\begin{taskbox}{T15 -- Build plain-language explanations}
\textbf{Goal.} Explain status using features, not vague claims.

\textbf{Steps.}
\begin{enumerate}
\item Start with z-score ranking: largest absolute personalized feature deviations become reasons.
\item For logistic regression, compute signed feature contributions if that model is used.
\item For random forest, use permutation/importances offline, but keep live explanations based on z-score and direction unless SHAP is fast enough.
\item Create restrained templates:
\begin{itemize}
\item ``Skin-response activity is above your quiet baseline.''
\item ``Pulse features are higher than your quiet baseline.''
\item ``Temperature trend is lower than your quiet baseline.''
\item ``Hand motion is high, so the model is paused.''
\end{itemize}
\item Never output medical explanations such as ``panic attack starting.''
\end{enumerate}

\begin{donebox}
Every colored status has 1--3 ranked reasons, every reason maps to a real feature, and no explanation
makes a diagnosis or unsupported prediction.
\end{donebox}
\end{taskbox}

\begin{taskbox}{T16 -- Build the FastAPI backend}
\textbf{Goal.} Serve live state to a dashboard without tying the app to hardware.

\textbf{Steps.}
\begin{enumerate}
\item Create \code{src/neuroshield/api/main.py}.
\item Implement endpoints:
\begin{itemize}
\item \code{GET /api/v1/health}
\item \code{GET /api/v1/system}
\item \code{POST /api/v1/session/start}
\item \code{POST /api/v1/calibration/start}
\item \code{GET /api/v1/status/latest}
\item \code{GET /api/v1/history}
\item \code{WS /ws/v1/live}
\end{itemize}
\item Add source mode options: \code{synthetic}, \code{replay}, and later \code{serial}.
\item Bind to \code{127.0.0.1} by default.
\item Add typed error codes for missing model, missing baseline, schema mismatch, stale source, and poor signal.
\end{enumerate}

\begin{donebox}
The backend starts with a replay file, health reports model/source readiness, WebSocket pushes status updates,
and tests cover success and error cases.
\end{donebox}
\end{taskbox}

\begin{taskbox}{T17 -- Build the dashboard}
\textbf{Goal.} Make the MVP usable from the first screen.

\textbf{Steps.}
\begin{enumerate}
\item Create the frontend under \code{app/} using Next.js or a simpler Streamlit dashboard if time is short.
\item Required visible sections: connection state, current status, quality row, reasons, latest values, short history, calibration controls.
\item Implement all states: disconnected, waiting, calibrating, green, amber, red, motion paused, poor signal, stale, and backend error.
\item Validate incoming backend messages in the browser. Unknown schema versions must show an error.
\item Add a history chart for the last few minutes of probability/status/features.
\end{enumerate}

\begin{donebox}
Running the replay backend and dashboard shows a complete live experience without hardware. The UI never shows an old green status as current after the backend stops.
\end{donebox}
\end{taskbox}

\begin{taskbox}{T18 -- Add dataset loader for one external validation dataset}
\textbf{Goal.} Prove the pipeline can handle data outside WESAD.

\textbf{Steps.}
\begin{enumerate}
\item Choose Stress-Predict first if fast validation is needed; choose Nurse Stress if testing messy real-world robustness is more important.
\item Follow the relevant download instructions from Section 2.
\item Build a loader that returns the same \code{SignalBundle}.
\item Document channel intersection: which features can and cannot be computed.
\item Run the frozen WESAD-trained model on compatible windows.
\item Save compatibility notes and metrics to \code{artifacts/metrics/external_validation_notes.md}.
\end{enumerate}

\begin{donebox}
At least one external dataset either runs through the frozen pipeline or has a written, specific incompatibility report.
\end{donebox}
\end{taskbox}

\begin{taskbox}{T19 -- Build the software acceptance replay}
\textbf{Goal.} Create the official proof that software is ready before hardware starts.

\textbf{Steps.}
\begin{enumerate}
\item Start from a clean terminal.
\item Run all backend tests and feature/model tests.
\item Generate or load the replay fixture.
\item Create a baseline from the quiet segment.
\item Start backend in replay mode.
\item Start dashboard.
\item Verify waiting, green, motion pause, recovery, amber/red, history, backend restart, stale state, and reconnect.
\item Save a screen recording or screenshots.
\item Write \code{artifacts/demo/software_acceptance.json} with commands, versions, commit, files, and pass/fail notes.
\end{enumerate}

\begin{donebox}
A person who did not write the code can follow the README and reproduce the full software-only demo.
This is the gate for starting hardware integration.
\end{donebox}
\end{taskbox}

\begin{taskbox}{T20 -- Prepare the hardware handoff contract}
\textbf{Goal.} Make later hardware work a source replacement, not a rewrite.

\textbf{Steps.}
\begin{enumerate}
\item Write \code{docs/hardware_handoff.md}.
\item Include the exact raw event schema, sample-rate targets, units, serial settings, health event requirements, and error handling.
\item Include a sample replay file and expected parser counters.
\item State that hardware acceptance begins only when the device emits \code{neuroshield.hw.v1} events accepted by the existing replay/serial adapter.
\end{enumerate}

\begin{donebox}
The future firmware builder can implement one job: emit the existing event contract. They do not need
to understand model training, dashboard internals, or feature extraction to begin.
\end{donebox}
\end{taskbox}

% ============================================================================
\section{Optional Tasks After Core MVP}

\begin{taskbox}{O1 -- Add PPG-DaLiA motion-quality analysis}
\textbf{Goal.} Replace guessed motion thresholds with evidence.

\textbf{Steps.}
\begin{enumerate}
\item Download PPG-DaLiA.
\item Load wrist PPG and acceleration plus chest ECG-derived ground truth.
\item Run pulse extraction under different motion levels.
\item Plot heart-rate error versus acceleration magnitude.
\item Pick conservative thresholds for \code{motion_paused}.
\end{enumerate}

\begin{donebox}
Motion thresholds are documented with a plot and a conservative chosen value.
\end{donebox}
\end{taskbox}

\begin{taskbox}{O2 -- Add SWELL-KW cognitive-load experiment}
\textbf{Goal.} Explore FocusGuard without changing the core M1 MVP.

\textbf{Steps.}
\begin{enumerate}
\item Request/download SWELL-KW if access is available.
\item Use only channels that map to the glove concept, mainly HRV/pulse-derived features and EDA.
\item Train a separate M2 classifier with subject-wise evaluation.
\item Label it experimental in the dashboard/report.
\end{enumerate}

\begin{donebox}
M2 has separate metrics and cannot be confused with the validated M1 model.
\end{donebox}
\end{taskbox}

\begin{taskbox}{O3 -- Add Fitbit trends screen}
\textbf{Goal.} Demonstrate long-term trend UI without claiming burnout prediction.

\textbf{Steps.}
\begin{enumerate}
\item Download the Kaggle Fitbit dataset.
\item Parse daily activity, minute-level heart rate, and sleep files.
\item Build rolling averages for resting heart-rate proxy, activity, and sleep duration.
\item Show trend cards in the dashboard using clear ``descriptive only'' wording.
\end{enumerate}

\begin{donebox}
The dashboard can display trend logic on real longitudinal data, and every label says descriptive trend rather than burnout detection.
\end{donebox}
\end{taskbox}

% ============================================================================
\section{Sprint Order}

\begin{longtable}{P{0.10\linewidth}P{0.22\linewidth}P{0.60\linewidth}}
\toprule
\textbf{Day} & \textbf{Primary focus} & \textbf{End-of-day evidence} \\
\midrule
1 & Scope, repo, Python environment & README, scope docs, \code{uv.lock}, package smoke test. \\
2 & WESAD download and loader & One WESAD subject plot and loader test. \\
3 & Signal bundle and features & Feature table for one subject, fixed \code{features-v1} columns. \\
4 & Labels and M1 training & LOSO metrics, dummy comparison, saved predictions. \\
5 & Model artifact and synthetic source & Versioned model, generated NDJSON replay fixture. \\
6 & Baseline, motion gate, live states & Replay produces green, motion pause, recovery, amber/red sequence. \\
7 & FastAPI backend & Health/status/WebSocket endpoints work with replay. \\
8 & Dashboard & Live UI shows status, reasons, quality, history, and stale/reconnect states. \\
9 & External validation and docs & Stress-Predict or Nurse loader/compatibility report, model card draft. \\
10 & Acceptance demo & Clean replay demo, tests, screenshots/recording, software acceptance JSON. \\
\bottomrule
\end{longtable}

% ============================================================================
\section{Final Software-First Definition of Done}

\begin{donebox}
The software-first MVP is complete only when all items below are true:
\begin{itemize}
\item WESAD is downloaded, documented, loadable, and not committed to Git.
\item A versioned \code{SignalBundle} abstraction exists.
\item \code{features-v1} has fixed columns, tests, and at least one WESAD output table.
\item M1 has subject-wise WESAD metrics, dummy comparison, saved predictions, and a model card.
\item A versioned model artifact and manifest load successfully.
\item Synthetic/replay NDJSON uses the same schema future hardware will use.
\item A personal baseline JSON is created, loaded, and applied to live features.
\item Motion/poor-signal windows cause explicit abstention.
\item The backend serves typed status over REST and WebSocket.
\item The dashboard shows live status, reasons, quality, calibration state, history, stale state, and reconnect.
\item At least one external validation dataset is attempted or has a documented access/incompatibility reason.
\item A clean-machine README can reproduce the software replay demo.
\item Hardware integration has not required changing the model, features, dashboard, or API contracts.
\end{itemize}
\end{donebox}

% ============================================================================
\section{Source Links}

\begin{itemize}
\item WESAD, UCI Machine Learning Repository:
\url{https://archive.ics.uci.edu/dataset/465/wesad+wearable+stress+and+affect+detection}
\item PPG-DaLiA, UCI Machine Learning Repository:
\url{https://archive.ics.uci.edu/dataset/495/ppg+dalia}
\item Stress-Predict Dataset, GitHub:
\url{https://github.com/italha-d/Stress-Predict-Dataset}
\item Nurse Stress Dataset, Dryad:
\url{https://datadryad.org/dataset/doi:10.5061/dryad.5hqbzkh6f}
\item SWELL-KW dataset page:
\url{http://cs.ru.nl/~skoldijk/SWELL-KW/Dataset.html}
\item Fitbit Fitness Tracker Data, Kaggle:
\url{https://www.kaggle.com/datasets/arashnic/fitbit}
\end{itemize}

\end{document}
